\section{Izomorfizmy, monomorfizmy i epimorfizmy}\label{sec:morfizmy}

\begin{definicja}[Izomorfizm]
Morfizm $f:A\to B$ w kategorii $\mathcal C$ jest \emph{izomorfizmem}, jeśli istnieje $g:B\to A$ takie, że $g\circ f=\id_A$ i $f\circ g=\id_B$. Morfizm $g$ nazywamy odwrotnością $f$ i oznaczamy $f^{-1}$. Piszemy $A\cong B$, jeśli istnieje izomorfizm między $A$ i $B$.
\end{definicja}

\begin{stwierdzenie}
Odwrotność izomorfizmu jest jednoznaczna. Złożenie dwóch izomorfizmów jest izomorfizmem, a odwrotnością złożenia jest złożenie odwrotności w odwrotnej kolejności.
\end{stwierdzenie}
\begin{proof}
Jeśli $g,h:B\to A$ są odwrotnościami $f:A\to B$, to $g=g\circ(f\circ h)=(g\circ f)\circ h=h$. Dla izomorfizmów $f:A\to B$, $u:B\to C$ mamy $(u\circ f)\circ(f^{-1}\circ u^{-1})=\id_C$ oraz $(f^{-1}\circ u^{-1})\circ(u\circ f)=\id_A$, korzystając z łączności.
\end{proof}

\begin{przyklad}[Co oznacza izomorfizm?]
W $\Set$ jest nim bijekcja, w $\Grp$ izomorfizm grup, w $\Vect_k$ odwracalne odwzorowanie liniowe, a w $\Top$ homeomorfizm. W kategorii pochodzącej z porządku częściowego jedynymi izomorfizmami są identyczności (bo $x\leq y$ i $y\leq x$ pociągają $x=y$). W kategorii praporządku różne obiekty mogą być izomorficzne.
\end{przyklad}

\begin{definicja}[Monomorfizm i epimorfizm]\label{def:mono-epi}
Morfizm $f:A\to B$ jest \emph{monomorfizmem}, jeśli dla każdego obiektu $X$ i wszystkich $u,v:X\to A$ zachodzi $f\circ u=f\circ v\Rightarrow u=v$. Morfizm $f:A\to B$ jest \emph{epimorfizmem}, jeśli dla każdego obiektu $Y$ i wszystkich $u,v:B\to Y$ zachodzi $u\circ f=v\circ f\Rightarrow u=v$.
\end{definicja}

Inaczej mówiąc, monomorfizmy można skracać z lewej strony równania, a epimorfizmy z prawej. Definicje te nie używają elementów ani obrazu odwzorowania.

\begin{twierdzenie}\label{tw:set-mono-epi}
Dla funkcji $f:A\to B$ w $\Set$ zachodzą równoważności: $f$ jest monomorfizmem wtedy i tylko wtedy, gdy jest różnowartościowa; $f$ jest epimorfizmem wtedy i tylko wtedy, gdy jest surjektywna.
\end{twierdzenie}
\begin{proof}
Jeżeli $f$ jest różnowartościowa, to równość $f\circ u=f\circ v$ daje $u(x)=v(x)$ dla każdego $x$. Odwrotnie, jeśli $f(a)=f(a')$, rozważmy funkcje $u,v:\{*\}\to A$ dane przez $u(*)=a$, $v(*)=a'$. Moniczność $f$ daje $u=v$, a więc $a=a'$.

Jeżeli $f$ jest surjektywna, to z $u\circ f=v\circ f$ wynika $u=v$ przez sprawdzenie na dowolnym $b=f(a)$. Jeśli $f$ nie jest surjektywna, wybierzmy $b_0\notin f(A)$ oraz funkcje $u,v:B\to\{0,1\}$, gdzie $u$ jest stale równa $0$, a $v(b_0)=1$ i $v(b)=0$ dla $b\neq b_0$. Mamy $u\neq v$, lecz $u\circ f=v\circ f$.
\end{proof}

\begin{przyklad}[Uwaga na kategorię]\label{przyklad:epi-pierscienie}
W kategorii $\Ring$ inkluzja $\mathbb Z\hookrightarrow\mathbb Q$ jest epimorfizmem, chociaż nie jest surjektywna. Istotnie, jeżeli $u,v:\mathbb Q\to R$ są homomorfizmami pierścieni z jedynką, które zgadzają się na $\mathbb Z$, to dla $n\in\mathbb Z\setminus\{0\}$ elementy $u(1/n)$ i $v(1/n)$ są obydwa odwrotnościami tego samego elementu $u(n)=v(n)$, a więc są równe. W konsekwencji $u(a/n)=v(a/n)$ dla wszystkich liczb wymiernych. Zatem $u=v$.

Przeciwnie, w kategorii pochodzącej z dowolnego porządku częściowego \emph{każdy istniejący} morfizm jest jednocześnie monomorfizmem i epimorfizmem, ponieważ każdy zbiór $\Hom(x,y)$ ma co najwyżej jeden element. Nie oznacza to, że każdy morfizm jest tam izomorfizmem.
\end{przyklad}

\begin{definicja}[Sekcja i retrakcja]
Dla morfizmu $f:A\to B$ morfizm $s:B\to A$ spełniający $f\circ s=\id_B$ nazywamy \emph{sekcją} $f$ (prawą odwrotnością $f$); sam $f$ jest wówczas \emph{rozszczepionym epimorfizmem}. Morfizm $r:B\to A$ spełniający $r\circ f=\id_A$ nazywamy \emph{retrakcją} $f$ (lewą odwrotnością $f$); sam $f$ jest wówczas \emph{rozszczepionym monomorfizmem}.
\end{definicja}

\begin{stwierdzenie}
Każdy rozszczepiony monomorfizm jest monomorfizmem, a każdy rozszczepiony epimorfizm jest epimorfizmem.
\end{stwierdzenie}
\begin{proof}
Jeżeli $r\circ f=\id_A$ i $f\circ u=f\circ v$, to $u=r\circ f\circ u=r\circ f\circ v=v$. Jeżeli $f\circ s=\id_B$ i $u\circ f=v\circ f$, to $u=u\circ f\circ s=v\circ f\circ s=v$.
\end{proof}

\begin{przyklad}
W $\Vect_k$ zanurzenie $i:k\to k^2$, $i(x)=(x,0)$, ma lewą odwrotność $p:k^2\to k$, $p(x,y)=x$, bo $p\circ i=\id_k$. Zatem $i$ jest rozszczepionym monomorfizmem, a $p$ rozszczepionym epimorfizmem. W $\Set$ funkcja $\varnothing\to\{*\}$ jest monomorfizmem, ale nie ma lewej odwrotności: nie istnieje funkcja $\{*\}\to\varnothing$.
\end{przyklad}

\begin{uwaga}
Morfizm będący jednocześnie monomorfizmem i epimorfizmem nazywamy \emph{bimorfizmem}. Nie musi on być izomorfizmem; świadczy o tym przykład $2\to6$ w kategorii porządku podzielności z przykładu~\ref{przyklad:poset}. Kategorie, w których każdy taki morfizm jest izomorfizmem, nazywa się \emph{zrównoważonymi} (ang. \emph{balanced}). Kategoria $\Set$ jest zrównoważona na mocy twierdzenia~\ref{tw:set-mono-epi}.
\end{uwaga}


\begin{przyklad}[Funktor zapominania nie musi zachowywać epimorfizmów]
Dla pierścieni przemiennych z jedynką funktor $U:\Ring\to\Set$ zapomina strukturę. Jak pokazano w przykładzie~\ref{przyklad:epi-pierscienie}, morfizm $\mathbb Z\hookrightarrow\mathbb Q$ jest epimorfizmem w $\Ring$. Jego obraz przez $U$ nie jest epimorfizmem w $\Set$, gdyż nie jest surjektywny. To konkretny kontrprzykład do błędnego przekonania, że każdy funktor zachowuje epimorfizmy.
\end{przyklad}

\begin{przyklad}[Monomorfizmy w kategorii grup]
Homomorfizm grup $f:G\to H$ jest monomorfizmem wtedy i tylko wtedy, gdy jest różnowartościowy. Jeden kierunek wynika ze skracania funkcji różnowartościowych. Jeśli $f(a)=f(b)$, weźmy dwa homomorfizmy $u,v:\mathbb Z\to G$ określone przez $u(1)=a$ i $v(1)=b$. Wówczas $f\circ u=f\circ v$, więc moniczność wymusza $u=v$, a stąd $a=b$. W ten sposób dla grup rolę jednoelementowego zbioru testowego z kategorii $\Set$ pełni grupa wolna na jednym generatorze.
\end{przyklad}

\subsection{Reguły skracania i odwrotności jednostronne}

\begin{stwierdzenie}\label{cat:prop:skracanie}
Złożenie monomorfizmów jest monomorfizmem, a złożenie epimorfizmów jest epimorfizmem. Jeśli $g\circ f$ jest monomorfizmem, to $f$ jest monomorfizmem; jeśli $g\circ f$ jest epimorfizmem, to $g$ jest epimorfizmem. Każdy izomorfizm jest bimorfizmem.
\end{stwierdzenie}
\begin{proof}
Niech $f:A\to B$, $g:B\to C$. Jeśli $g\circ f\circ u=g\circ f\circ v$, moniczność $g$ daje $f\circ u=f\circ v$, a moniczność $f$ daje $u=v$. Jeśli $u\circ g\circ f=v\circ g\circ f$, epiczność $f$ daje $u\circ g=v\circ g$, a epiczność $g$ daje $u=v$.
Gdy $f\circ u=f\circ v$, po złożeniu z $g$ z lewej strony i skróceniu monomorfizmu $g\circ f$ dostajemy $u=v$. Analogicznie z $u\circ g=v\circ g$ i epiczności $g\circ f$ otrzymujemy $u=v$. Izomorfizm ma obie odwrotności jednostronne, zatem jest zarazem rozszczepionym monomorfizmem i rozszczepionym epimorfizmem.
\end{proof}

\begin{stwierdzenie}\label{cat:prop:bimorfizm-rozszczepiony}
Monomorfizm mający prawą odwrotność jest izomorfizmem. Epimorfizm mający lewą odwrotność jest izomorfizmem.
\end{stwierdzenie}
\begin{proof}
Jeśli $f\circ s=\id_B$ i $f$ jest monomorfizmem, to $f\circ(s\circ f)=f=f\circ\id_A$. Po skróceniu dostajemy $s\circ f=\id_A$. Jeśli $r\circ f=\id_A$ i $f$ jest epimorfizmem, to $(f\circ r)\circ f=f=\id_B\circ f$, więc $f\circ r=\id_B$.
\end{proof}

\subsection{Przykłady z algebry i topologii}

\begin{twierdzenie}[Epimorfizmy grup]\label{cat:thm:epi-grp}
Epimorfizmy w $\Grp$ to dokładnie homomorfizmy surjektywne.
\end{twierdzenie}
\begin{proof}
Surjektywność wystarcza, ponieważ funkcje o tej samej dziedzinie, zgadzające się na całym obrazie surjekcji, są równe. Niech teraz $f:G\to H$ nie będzie surjekcją i niech $K=f(G)$ będzie właściwą podgrupą $H$. Rozważmy zbiór $X=H\times\{0,1\}$ i homomorfizm $u:H\to\operatorname{Sym}(X)$ określony przez $u(h)(x,i)=(hx,i)$. Zdefiniujmy permutację $\tau$ zbioru $X$, która zamienia $(x,0)$ z $(x,1)$ dla $x\in K$, a wszystkie pozostałe punkty pozostawia na miejscu. Niech $v(h)=\tau\circ u(h)\circ\tau^{-1}$. Jest to homomorfizm, ponieważ sprzężenie przez ustaloną permutację zachowuje składanie.

Dla $k\in K$ zachodzi $kx\in K$ wtedy i tylko wtedy, gdy $x\in K$. Stąd $u(k)$ komutuje z $\tau$, więc $v(k)=u(k)$ i $v\circ f=u\circ f$. Jeżeli $h\in H\setminus K$, to jednak $u(h)(1,0)=(h,0)$, podczas gdy $v(h)(1,0)=\tau(h,1)=(h,1)$. Zatem $u\ne v$ i $f$ nie jest epimorfizmem.
\end{proof}

\begin{twierdzenie}[Morfizmy modułów]\label{cat:thm:mono-epi-mod}
W kategorii lewych unitarnych modułów nad pierścieniem $R$ z jedynką monomorfizmy są dokładnie homomorfizmami różnowartościowymi, a epimorfizmy dokładnie homomorfizmami surjektywnymi. Kategoria ta jest zrównoważona.
\end{twierdzenie}
\begin{proof}
Dla $f:M\to N$ różnowartościowość i surjektywność implikują odpowiednie skracanie funkcji. Jeśli $f(x)=f(y)$, dwa homomorfizmy $a,b:R\to M$ dane przez $a(r)=rx$, $b(r)=ry$ mają równe złożenia z $f$. Moniczność daje $a=b$, więc $x=a(1)=b(1)=y$. Jeśli $f$ nie jest surjektywna, moduł $N/f(M)$ jest niezerowy. Rzut ilorazowy $q:N\to N/f(M)$ różni się od homomorfizmu zerowego, ale $q\circ f=0$. Zatem $f$ nie jest epimorfizmem. Bijektywny homomorfizm modułów ma liniową odwrotność: dla $n=f(x)$, $n'=f(y)$ zachodzi $f^{-1}(rn+sn')=rx+sy$. Wykorzystujemy konstrukcje ilorazów i homomorfizmów z rozdziału 16 skryptu \emph{Algebra abstrakcyjna}.
\end{proof}

\begin{przyklad}[Bimorfizm bez rozszczepienia]\label{cat:ex:mono-epi-top}
W $\Top$ monomorfizmy to dokładnie ciągłe funkcje różnowartościowe: do dowodu konieczności używamy dwóch odwzorowań z przestrzeni jednopunktowej. Epimorfizmy to dokładnie ciągłe surjekcje. Jeśli ciągła funkcja $f:X\to Y$ nie jest surjekcją, wybieramy $y_0\in Y\setminus f(X)$ i dwie funkcje do zbioru $\{0,1\}$ różniące się jedynie w $y_0$, tak jak w dowodzie twierdzenia~\ref{tw:set-mono-epi}. Nadajemy zbiorowi docelowemu topologię antydyskretną $\{\varnothing,\{0,1\}\}$, dzięki czemu obie funkcje są ciągłe. Dowodzi to braku epiczności.

Identyczność ze zbioru $\{0,1\}$ z topologią dyskretną do tego samego zbioru z topologią antydyskretną jest ciągłą bijekcją, ale jej odwrotność nie jest ciągła (przeciwobraz singletonu nie jest otwarty). Jest zatem bimorfizmem, który nie jest izomorfizmem. Na mocy stwierdzenia~\ref{cat:prop:bimorfizm-rozszczepiony} nie jest też rozszczepiony z żadnej strony.
\end{przyklad}

\begin{przyklad}[Gładka bijekcja nie musi być dyfeomorfizmem]
Odwzorowanie $f:\mathbb R\to\mathbb R$, $f(x)=x^3$, jest gładką bijekcją, więc można je skracać z obu stron w kategorii rozmaitości. Odwrotność nie jest różniczkowalna w zerze, gdyż dla $t>0$ iloraz różnicowy wynosi $t^{1/3}/t=t^{-2/3}$ i jest nieograniczony, gdy $t\to0^+$. Stąd $f$ nie jest izomorfizmem w tej kategorii. Pojęcie dyfeomorfizmu pochodzi z rozdziału 1 skryptu \emph{Geometria i topologia różniczkowa}.
\end{przyklad}

\begin{przyklad}[Lokalizacja jako epimorfizm]\label{cat:ex:epi-lokalizacja}
Uogólnijmy przykład~\ref{przyklad:epi-pierscienie}. Niech $R$ będzie pierścieniem przemiennym z jedynką, $S$ jego zbiorem multiplikatywnym, a $j:R\to S^{-1}R$ kanonicznym homomorfizmem. Własność uniwersalna lokalizacji (rozdział 12 skryptu \emph{Algebra abstrakcyjna}) mówi, że homomorfizm z $R$, który przekształca elementy $S$ w elementy odwracalne, ma dokładnie jedno przedłużenie wzdłuż $j$. Jeśli $u\circ j=v\circ j$, wspólne złożenie spełnia ten warunek, ponieważ $j(s)$ jest odwracalny dla każdego $s\in S$. Jednoznaczność daje $u=v$. Zatem $j$ jest epimorfizmem w $\Ring$ bez założenia, że $j$ jest różnowartościowy.
\end{przyklad}

\subsection{Monomorfizmy i epimorfizmy regularne}

\begin{definicja}[Regularność morfizmu]\label{cat:def:regularne}
Monomorfizm regularny to morfizm $m:A\to B$, dla którego istnieją $u,v:B\to C$ takie, że $u\circ m=v\circ m$, a każdy $t:X\to B$ spełniający $u\circ t=v\circ t$ ma dokładnie jedną faktoryzację $t=m\circ w$ z $w:X\to A$. Dualnie epimorfizm regularny $e:B\to A$ spełnia $e\circ u=e\circ v$ dla pewnych $u,v:C\to B$, a każdy $t:B\to X$ spełniający $t\circ u=t\circ v$ ma dokładnie jedną faktoryzację $t=w\circ e$. Są to konstrukcje uniwersalne omawiane w rozdziale~\ref{sec:equalizatory}.
\end{definicja}

\begin{stwierdzenie}\label{cat:prop:regularne-mono-epi}
Monomorfizm regularny jest monomorfizmem, a epimorfizm regularny jest epimorfizmem.
\end{stwierdzenie}
\begin{proof}
Jeśli $m\circ a=m\circ b=t$, to $t$ spełnia $u\circ t=v\circ t$ i zarówno $a$, jak i $b$ są jego faktoryzacjami. Jednoznaczność daje $a=b$. Jeśli $a\circ e=b\circ e=t$, to $t\circ u=t\circ v$ i obie faktoryzacje $a,b$ muszą być równe.
\end{proof}

\begin{przyklad}
W $\Set$ każda iniekcja $m:A\to B$ jest regularna: niech $u:B\to\{0,1\}$ będzie stale zerowa, a $v$ będzie zerowa dokładnie na $m(A)$ i jedynkowa poza nim. Funkcja $t:X\to B$ spełnia $u\circ t=v\circ t$ dokładnie wtedy, gdy $t(X)\subseteq m(A)$; wtedy $w(x)$ jest jedynym elementem $A$ o obrazie $t(x)$. Każda surjekcja $e:B\to A$ jest regularna: bierzemy relację $C=\{(b,b')\in B^2:e(b)=e(b')\}$ oraz jej dwa rzuty $u,v:C\to B$. Jeśli $t\circ u=t\circ v$, to $t$ jest stała na włóknach $e$. Wzór $w(a)=t(b)$ dla dowolnego $b$ z $e(b)=a$ jest poprawny i jednoznacznie określa $w$; $t=w\circ e$. Nie trzeba tutaj wybierać globalnego przekroju surjekcji.
\end{przyklad}

